\documentclass{article}

\usepackage{amsmath}
\usepackage{amsthm}
\usepackage{amssymb}

\title{Basic Geometry}
\author{Rovenszky Robin}
\date{Last updated: 10 March 2026}

\newtheorem{theorem}{Theorem}

\begin{document}

\maketitle

\begin{theorem}[Pythagorean Theorem]
In a right triangle with legs of lengths $a$ and $b$, and hypotenuse of length $c$, the following relation holds:
\[
a^2 + b^2 = c^2.
\]
\end{theorem}

\begin{proof}
Consider a square with side length $(a+b)$. Inside the square, place four identical right triangles with legs $a$ and $b$ and hypotenuse $c$, arranged so that their hypotenuses form a smaller square in the center.

The area of the large square is
\[
(a+b)^2.
\]

This area is also equal to the sum of the areas of the four triangles and the small central square.

Each triangle has area
\[
\frac{1}{2}ab,
\]
so the four triangles together have area
\[
4 \cdot \frac{1}{2}ab = 2ab.
\]

The remaining central square has side length $c$, so its area is
\[
c^2.
\]

Thus,
\[
(a+b)^2 = 2ab + c^2.
\]

Expanding the left-hand side gives
\[
a^2 + 2ab + b^2 = 2ab + c^2.
\]

Subtracting $2ab$ from both sides yields
\[
a^2 + b^2 = c^2.
\]

This proves the theorem.
\end{proof}

\newpage

\section{Related Problems}
\subsection{Problem 1.}

First Pythagorean Theorem:
\begin{equation}
    12^2 + x^2 = 15^2
\end{equation}
\begin{equation}
    144+x^2 = 225
\end{equation}
\begin{equation}
    x^2 = 81
\end{equation}
\begin{equation}
    x = 9
\end{equation}
Second Pythagorean Theorem:
\begin{equation}
    12^2 + y^2 = 13^2
\end{equation}
\begin{equation}
    144 + y^2 = 169
\end{equation}
\begin{equation}
    y^2 = 25
\end{equation}
\begin{equation}
    y = 5
\end{equation}
Thus our final answer is $x + y = 14$. \qed

\subsection{Problem 2.}

First Pythagorean Theorem:
\begin{equation}
    x^2 + (12-r)^2 = (12+r)^2
\end{equation}
Second Pythagorean Theorem:
\begin{equation}
    x^2 + r^2 = (24-r)^2
\end{equation}
Evaluating (9) - (10):
\begin{equation}
    (12-r)^2 - r^2 = (12+r)^2 - (24-r)^2
\end{equation}
\begin{equation}
    (12-2r) \cdot 12 = 36 \cdot (2r - 12)
\end{equation}
Dividing by 12 yields:
\begin{equation}
    (12-2r) = (2r - 12) \cdot 3
\end{equation}
Then dividing by 2:
\begin{equation}
    6-r = 3 \cdot (r-6)
\end{equation}
\begin{equation}
    6-r = 3r - 18
\end{equation}
Adding 18+r:
\begin{equation}
    24 = 4r
\end{equation}
\begin{equation}
    r = 6
\end{equation}
\qed

\subsection{Problem 3.}
\begin{theorem}[Theorem 1.]
For a parallelogram, we have
\begin{equation}
    e^2 + f^2 = a^2 \cdot2 + b^2 \cdot2
\end{equation}
\begin{proof}
First Pythagorean Theorem:
\begin{equation}
    (a-x)^2 + m^2 = e^2
\end{equation}
Second Pythagorean Theorem:
\begin{equation}
    (a+x)^2 + m^2 = f2
\end{equation}
Evaulating (19) + (20):
\begin{equation}
    a^2 - 2ax + x^2 + m^2 + a^2 + 2ax + x^2 + m^2 = e^2 + f^2
\end{equation}
Simplifying and invoking $m^2 + x^2 = b^2$:
\begin{equation}
    2a^2 + 2b^2 = e^2 + f^2
\end{equation}
\end{proof}

\end{theorem}

\subsection{Problem 4.}
Let the small circle's radius be $r$.\\
First Pythagorean Theorem:
\begin{equation}
    x^2 + (12 - r)^2 = (12 + r)^2
\end{equation}
Second Pythagorean Theorem:
\begin{equation}
    r^2 + x^2 = (24-r)^2
\end{equation}
See (1.2) Related Problems: Problem 2. for solution. These are two equivalent problems, under a 90 degree rotation. See figures for clarification.

\subsection{Problem 5.}
Let's put the circles inside rectangles instead of other circles!\\
We seek to find the height of the rectangle given the diameter of the small circles.\\
$r_0 = 12$

\subsection{Problem 6.}
Another circle-inside-rectangle problem. The rectangle is 60x48, what is the radius of the small circles?

\section{Other}
\subsection{What is $\sqrt{2}$?}
We understand that it must exist, since in a right trinagle with legs $a = 1, b= 1, \alpha = 90$, where $\alpha$ is their enclosed angle, the hypothenuse has length $\sqrt{2}$. So what is $\sqrt{2}$?
\\
\begin{theorem}
    For an x satisfying $x^2 = 2$, x may also satisfy $x = \frac{p}{q}$ with $p, q \in \mathbb{Z}$, $gcd(p, q) = 1$
\end{theorem}
\begin{proof}
    Assume the statement is true. Then
    \begin{equation}
        \frac{p^2}{q^2} = 2 \implies p^2 = 2q^2 \implies 2 \mid p^2 \implies 2 \mid p \implies 4 \mid p^2 \implies
    \end{equation}
    \begin{equation}
        \implies 4 \mid 2q^2 \implies 2 \mid q^2 \implies 2 \mid q \implies \gcd(p, q) \ge 2 \Rightarrow\!\Leftarrow
    \end{equation}
    Or alternatively, we could have also finished the proof via
    $p^2 = 2q^2$ a 2 kitevője balra ps, jobbra ptlan
\end{proof}

\end{document}